% ============================================================================
% Fichier  : partie1/P01_C03_pdca.tex
% Titre    : Le Cycle PDCA et l'Amélioration Continue
% Auteur   : MINKA MI NGUIDJOÏ Thierry Emmanuel
% ============================================================================
% Sources : ISO/IEC 27001:2022 (approche processus). Mairie d'Eséka : cas fictif.
% ============================================================================

\chapter{Le Cycle PDCA et l'Amélioration Continue}
\label{chap:pdca}
\chaptermark{Le cycle PDCA}

\epigraph{%
    « Un système de sécurité qui ne tourne pas se dégrade~:\\
    ce qui était sûr hier ne l'est plus demain. »%
}{Principe du cours}

\niveauChapitre{Fondamental}{Les Chapitres~\ref{chap:cia} et~\ref{chap:enjeux}.}

\begin{boxobjectifs}
\begin{itemize}
    \item Énoncer les quatre étapes du cycle \sigle{PDCA}.
    \item Comprendre pourquoi la sécurité est un cycle et non un projet.
    \item Relier chaque étape aux parties du cours et aux phases du fil rouge.
    \item Distinguer correction, action corrective et amélioration.
    \item Situer l'amélioration continue comme exigence de la norme.
\end{itemize}
\end{boxobjectifs}

\bigskip

% ============================================================================
\section*{Question directrice : pourquoi la sécurité n'est-elle jamais « finie » ?}
% ============================================================================

La mairie pourrait croire qu'après avoir installé des mesures, le travail est
accompli. Ce chapitre montre que la sécurité se dégrade si on ne l'entretient
pas, et que le cycle \sigle{PDCA} est l'outil de cet entretien.

\medskip
\textbf{Plan du chapitre.}
\begin{enumerate}
    \item Les quatre étapes (section~\ref{sec:c03-etapes}).
    \item Un cycle, pas un projet (section~\ref{sec:c03-cycle}).
    \item PDCA et plan du cours (section~\ref{sec:c03-plan}).
    \item Corriger, améliorer (section~\ref{sec:c03-corriger}).
    \item L'amélioration continue, une exigence (section~\ref{sec:c03-exigence}).
\end{enumerate}

\separateur

% ============================================================================
\section{Les Quatre Étapes}
\label{sec:c03-etapes}
% ============================================================================

\begin{boxdef}
Le \termecle{cycle PDCA} (\emph{Plan, Do, Check, Act}), ou roue de Deming,
organise l'amélioration continue en quatre étapes répétées~:
\termecle{planifier}, \termecle{déployer}, \termecle{contrôler}, \termecle{agir}.
\end{boxdef}

\begin{tableaucours}{L{2.6cm}L{4.4cm}L{5.2cm}}{%
    \tableauHeader{Étape} &
    \tableauHeader{But} &
    \tableauHeader{À la mairie}
}
Planifier & Fixer objectifs et mesures & Apprécier les risques, écrire la politique \\
Déployer & Mettre en \oe uvre & Déployer les mesures, former les agents \\
Contrôler & Vérifier l'efficacité & Mesurer, auditer, revue de direction \\
Agir & Améliorer & Corriger les écarts, décider le cycle suivant \\
\end{tableaucours}
\vspace{-0.8em}
\begin{center}{\small\itshape Tableau~3.1. Le cycle PDCA, étape par étape}\end{center}

% ============================================================================
\section{Un Cycle, pas un Projet}
\label{sec:c03-cycle}
% ============================================================================

Un projet a une fin~; un cycle recommence. La menace évolue, les agents
changent, les systèmes vieillissent~: une mesure efficace en janvier peut ne
plus l'être en décembre. Le \sigle{PDCA} rattrape cette dérive en imposant de
revenir régulièrement sur ce qui a été décidé.

\begin{boxscenario}{Opération NYONG. Le piège du « c'est fait »}
\textbf{Constatation.} Six mois après avoir imposé des mots de passe
individuels, la mairie découvre que trois agents les ont de nouveau partagés.

\medskip
\textbf{Analyse} (confiance élevée). La mesure avait été déployée (\emph{Do})
mais jamais contrôlée (\emph{Check}). Sans l'étape de contrôle, une mesure se
relâche. Le cycle n'est pas une formalité~: c'est ce qui empêche le retour en
arrière.
\end{boxscenario}

% ============================================================================
\section{PDCA et Plan du Cours}
\label{sec:c03-plan}
% ============================================================================

Le cours épouse le cycle. Les parties II et III \emph{planifient} (gouvernance,
politique, cadre \norme{ISO/IEC 27001})~; la partie IV \emph{planifie} par le
risque~; la partie V \emph{déploie} les mesures~; la partie VI \emph{contrôle}
et \emph{agit} (audit, certification, incidents, amélioration).

\begin{boxmetier}{RSSI}
\textbf{Tenir la roue.} Le \sigle{RSSI} n'a pas pour mission de « finir » la
sécurité, mais de faire tourner la roue~: chaque année, replanifier au vu des
nouveaux risques et des écarts constatés. Son tableau de bord suit le cycle.
\end{boxmetier}

% ============================================================================
\section{Corriger, Améliorer}
\label{sec:c03-corriger}
% ============================================================================

\begin{boxdef}
\termecle{Correction}~: action immédiate qui traite un écart constaté (rétablir
un service). \termecle{Action corrective}~: action qui traite la \emph{cause}
de l'écart pour qu'il ne se reproduise pas. \termecle{Amélioration}~:
renforcement au-delà de la simple conformité.
\end{boxdef}

Rétablir les postes après la clé \texttt{USB} était une \emph{correction}.
Interdire les clés personnelles et imposer des sauvegardes est une \emph{action
corrective}. Automatiser la détection des clés non autorisées serait une
\emph{amélioration}. La norme attend les trois, dans cet ordre de maturité.

% ============================================================================
\section{L'Amélioration Continue, une Exigence}
\label{sec:c03-exigence}
% ============================================================================

L'amélioration continue n'est pas une bonne intention~: c'est la clause~10 de
la norme \norme{ISO/IEC 27001} (Chapitre~\ref{chap:clauses710}). Un
\sigle{SMSI} qui ne s'améliore pas n'est pas conforme. Le \sigle{PDCA} est donc
à la fois une philosophie et une exigence vérifiée en audit.

\begin{boxinfo}
\textbf{À retenir pour tout le cours.} Chaque fois qu'une partie s'ouvre,
demandez-vous à quelle étape du \sigle{PDCA} elle correspond. La structure du
cours \emph{est} le cycle.
\end{boxinfo}

\separateur

\begin{boxresume}
\textbf{Ce qu'il faut retenir de ce chapitre~:}
\begin{itemize}
    \item Le \textbf{PDCA} (planifier, déployer, contrôler, agir) organise
          l'amélioration continue.
    \item La sécurité est un \textbf{cycle}, pas un projet~: sans entretien,
          elle se dégrade.
    \item Le \textbf{cours} suit les étapes du cycle.
    \item \textbf{Correction}, \textbf{action corrective} et
          \textbf{amélioration} se distinguent par ce qu'elles traitent.
    \item L'\textbf{amélioration continue} est une exigence de la norme
          (clause~10).
\end{itemize}
\end{boxresume}

\begin{boxinfo}
\textbf{Points de contrôle de l'élève.} Je sais~: énoncer les quatre étapes~;
expliquer pourquoi la sécurité est un cycle~; relier le PDCA au plan du cours~;
distinguer correction, action corrective et amélioration~; rattacher
l'amélioration continue à la norme.
\end{boxinfo}

\bigskip

\begin{boxexo}[title={\ding{45}\quad Exercice~3.1. Les quatre étapes \niveaubase}]
Nommez les quatre étapes du \sigle{PDCA} et donnez une action de la mairie pour
chacune.
\end{boxexo}

\begin{boxexo}[title={\ding{45}\quad Exercice~3.2. Cycle ou projet \niveaubase}]
Expliquez en deux phrases pourquoi la sécurité n'est pas un projet qui se
termine.
\end{boxexo}

\begin{boxexo}[title={\ding{45}\quad Exercice~3.3. Trois mots \niveaubase}]
Distinguez correction, action corrective et amélioration par un exemple de
chacune.
\end{boxexo}

\begin{boxexo}[title={\ding{45}\quad Exercice~3.4. L'étape manquante \niveaumoyen}]
Dans le cas des mots de passe repartagés, quelle étape du cycle a manqué~?
Justifiez.
\end{boxexo}

\begin{boxexo}[title={\ding{45}\quad Exercice~3.5. Rattacher au cycle \niveaumoyen}]
À quelle étape du PDCA rattachez-vous~: un audit interne~; une formation~; une
appréciation des risques~; une revue de direction~?
\end{boxexo}

\begin{boxexo}[title={\ding{45}\quad Exercice~3.6. De la correction à la cause \niveaumoyen}]
Pour l'incident de la clé \texttt{USB}, proposez une correction, une action
corrective et une amélioration.
\end{boxexo}

\begin{boxexo}[title={\ding{45}\quad Exercice~3.7. Dérive \niveaumoyen}]
Donnez deux raisons pour lesquelles une mesure efficace peut cesser de l'être
avec le temps.
\end{boxexo}

\begin{boxexo}[title={\ding{45}\quad Exercice~3.8. Exigence \niveaumoyen}]
Pourquoi un SMSI qui ne s'améliore jamais n'est-il pas conforme à
\norme{ISO/IEC 27001}~?
\end{boxexo}

\begin{boxexo}[title={\ding{45}\quad Exercice~3.9. Planifier un tour \niveauexpert}]
Décrivez un cycle PDCA annuel complet pour la mairie d'Eséka, de la
planification à l'amélioration.
\end{boxexo}

\begin{boxexo}[title={\ding{45}\quad Exercice~3.10. Note au conseil \niveauexpert}]
Rédigez une note expliquant au conseil municipal pourquoi la sécurité
nécessite un budget récurrent, et non une dépense unique.
\end{boxexo}

\bigskip

\begin{boxinfo}
\textbf{Pour aller plus loin.}
\begin{itemize}
    \item ISO/IEC 27001, approche processus et clause~10~\cite{iso27001}.
\end{itemize}
\end{boxinfo}

\begin{boxinfo}
\textbf{Transition.} Faire tourner la roue suppose une organisation qui la
porte~: le Chapitre~\ref{chap:gouvernance} ouvre la gouvernance de la sécurité.
\end{boxinfo}
